% =====================================================================
%   4. AgentBench 的评测 (Evaluation)
% =====================================================================

\section{AgentBench 的评测}

我们对 29 个 LLM（涵盖 API 商用模型与开源 LLM）进行了大规模评测，以期
系统呈现 LLM-as-Agent 的当前性能现状。同时，本文设计并开源了一套
"即插即用"的评估工具包，便于开展相关的 LLM-as-Agent 研究。

\begin{table}[t]
    \centering
    \small
    \caption{
        \textbf{表 2：AgentBench 中 8 个环境的数据统计与评测指标}。
        "SR" 表示成功率（Success Rate）；"\#平均轮数" 为单条任务所需的
        估计交互轮数；在 "\#Dev" 与 "\#Test" 中给出查询样本数与总估计
        交互轮数。此外，"Weight\textsuperscript{$-1$}" 表示在所有模型上
        该任务的平均得分。详见正文~\ref{sec:eval-setup}~节与附录 B--I。
    }
    \label{tab:stats}
    \begin{tabularx}{\textwidth}{l *{8}{>{\centering\arraybackslash}X}}
        \toprule
        & \textbf{操作系统} & \textbf{数据库} & \textbf{知识图谱}
        & \textbf{数字卡牌} & \textbf{横向谜题} & \textbf{家务整理}
        & \textbf{网页购物} & \textbf{网页浏览}\\
        \midrule
        \#平均轮数        & 8   & 5   & 15  & 30  & 25  & 35  & 5   & 10 \\
        评测指标         & SR  & SR  & F1  & 奖励 & 游戏进度 & SR  & 奖励 & 步骤 SR \\
        \#Dev            & 26/240   & 60/300    & 20/300   & 12/360    & 20/500    & 20/700   & 80/400   & 31/400 \\
        \#Test           & 144/1200 & 300/1500  & 150/2250 & 20/600    & 50/1250   & 50/1750  & 200/1000 & 100/1000 \\
        Weight\textsuperscript{$-1$} & 10.8 & 13.0 & 13.9 & 12.0 & 3.5 & 13.0 & 30.7 & 11.6 \\
        \bottomrule
    \end{tabularx}
\end{table}

\subsection{评测设置}

\paragraph{数据集统计。}
AgentBench 各数据集的统计信息见表~\ref{tab:stats}。为简洁起见，后文将
使用各数据集的缩写形式。所有数据集均为实用的多轮交互任务，其单条问题的
预计求解轮数介于 5 到 50 之间。我们为每个数据集提供开发集（Dev）与
测试集（Test）两个划分，且所有数据集均已开源。

在设计 AgentBench 时，我们细致地平衡了评估的全面性与效率，因为 LLM 的
多轮交互可能非常耗时。我们将开发集与测试集的样本量分别设置为 269 与
1\,014，对应的推理调用次数约为 3\,000 与 11\,000——这与 MMLU
\citep{hendrycks2021b} 所需的推理调用量相当。

\paragraph{参与评测的 LLM。}
作为对现有 LLM 在 LLM-as-Agent 任务上的首次系统性评测尝试，本文共纳入
29 个模型进行评估，可大致划分为以下两类：
\begin{itemize}
    \item \textbf{API 类商用 LLM}：主要为参数规模未公开的 LLM API
          （参见表~\ref{tab:models}）。由于投入巨大，其性能通常更佳；
    \item \textbf{开源（OSS）LLM}：主要来自学术界与部分企业
          （参见表~\ref{tab:models}）。受算力资源所限，本文仅纳入
          规模不超过 70B 的 OSS LLM。需要指出的是，该规模已涵盖了
          绝大多数在对话或指令任务上经过特定微调、性能优异的开源 LLM
          （仅有极少数例外）。
\end{itemize}

\paragraph{工具包：基于 API 中心化思想与环境隔离以简化 LLM 评测。}
随着 LLM 系统在复杂度上不断提升，且主要通过 API 提供访问，我们开发了一套
与 API 化理念相契合的评估工具包。该工具包专为与 API 交互而精心设计，
简化了对不同 LLM 进行适配与测试的流程。研究者若希望在 AgentBench 上评测
其 LLM，仅需部署一个可通过 HTTP 协议访问的模型服务即可。

此外，处理多样且复杂的交互环境本身就是一项重大挑战。若试图对所有环境进行
统一配置，往往困难重重并易引发冲突。为此，我们采用了两项关键策略：
\textbf{其一}，我们将含复杂环境的任务封装为 Docker 镜像，研究者只需挂载代码路径
并启动评估即可便捷使用这些镜像；\textbf{其二}，我们将每项任务进一步拆分为
独立的 worker，从而保证各任务环境的隔离，规避冲突。（详见附录~\ref{app:framework}）

\paragraph{评估提示词设置。}
为兼顾当前大多数对话模型，本文采用双角色——\code{user}（即指令与环境反馈）
与 \code{agent}——交替发言的对话范式。我们将交互轨迹记录为一段对话历史
$(u_0, a_0, \dots, u_k, a_k)$，其中 $u_i$、$a_i$ 分别表示第 $i$ 轮
对话历史中用户与智能体的发言。在执行推理时，对话历史遵循
$(u_0, a_0, \dots, u_k)$ 的格式。我们选取最小的 $r$，使得
$(u_0, a_r, u_{r+1}, \dots, u_k)$ 中所有 token 的总数
\footnote{由于各模型的分词器不同，我们简单地按以下规则估算 token 数：
长度为 $n$ 的单词占据 $\lceil n/6 \rceil$ 个 token，非空白字符则各占 1 个 token。}
不超过 3\,500；随后在 $u_0$ 中追加一句
\code{"[NOTICE] $2r$ messages are omitted."}。最终，
$(u_0, a_r, u_{r+1}, \dots, u_k)$ 被视为多轮对话格式下的最终输入。

不过，为兼顾非对话模型，我们额外引入了后处理器：对支持多轮的对话模型，
我们直接以对话历史作为输入；对仅支持文本补全的模型（例如
\code{text-davinci-003}），则在历史的每条内容前添加前缀
\code{"USER:"} 或 \code{"AGENT:"}，并在末尾追加字符串
\code{"AGENT:"}，以使模型生成智能体的回复。

在任务提示词的构造上，我们参考了 \citet{yao2023b} 的格式，将"思维"
（Thought，用于 CoT）与"动作"（Action）整合在单轮回复中。通常，任务指令
中会附带一段简单的 CoT 演示示例，以引导模型生成符合规范的输出。为保证
结果可复现，所有任务的推理均采用 \code{temperature=0}（即贪心解码），
与 \citet{wei2022b} 的做法保持一致。

\paragraph{总分计算方式。}
我们观察到，各任务的得分分布存在显著差异，因为不同任务的难度水平不一。
若直接对各任务的得分取算术平均，将被那些普遍得分偏高的任务（例如
我们观察到的网页购物）所主导，进而掩盖得分偏低的任务，难以契合
AgentBench 的评测目的。

因此，本文计算总分时，首先将每个任务在所有评估模型上的平均得分统一
归一化至 1，然后对每个模型在各任务上的得分取平均（参见表~\ref{tab:stats}）。
为便于后续研究在标准化与简化方面复用，本文以全部被测 LLM 在每项任务上的
平均得分的倒数，作为后续总分计算中的固定权重；总分即每项任务得分与其对应
权重乘积的平均值。这一方式保证了评估的公平性与一致性，并便于后续研究
进行对比与分析。

\subsection{主要结果}

表~\ref{tab:main-results} 报告了 AgentBench 的总体得分与各数据集的具体得分。
令人惊讶的是，在这一具有挑战性的基准上，我们发现部分顶级 LLM 已具备应对
真实世界环境交互的扎实能力：例如 \code{gpt-4} 在 AgentBench 8 个数据集中的
6 个上均取得最佳表现；在家务整理任务上，其成功率高达 78\%，表明该场景下
已具备一定的实用性。\code{claude-2} 与 \code{claude} 紧随其后，且均显著
优于 \code{gpt-3.5-turbo}。尽管其它 API 类 LLM 的表现相对逊色，但无论在哪
些任务上，它们中的大多数都能解决相当比例的问题。所有 API 类 LLM 在
AgentBench 上的总体得分均超过 1.00。

\begin{table*}[t]
    \centering
    \small
    \caption{
        \textbf{表 3：AgentBench 测试集（标准）结果}。
        顶级商用 LLM（如 \code{gpt-4}）与 OSS LLM 竞品之间存在明显的
        性能差距。"VER" 表示模型版本；"OA" 表示 AgentBench 总体得分，
        即所有环境得分的加权平均（参见~\ref{sec:eval-setup}~节）。
    }
    \label{tab:main-results}
    \begin{tabular}{ll c c *{4}{c} *{3}{c}}
        \toprule
        \multirow{2}{*}{\textbf{LLM 类型}} & \multirow{2}{*}{\textbf{模型}} & \multirow{2}{*}{\textbf{版本}}
        & \multirow{2}{*}{\textbf{OA}}
        & \multicolumn{4}{c}{\textbf{代码接地}} & \multicolumn{3}{c}{\textbf{游戏接地}} \\
        \cmidrule(lr){5-8} \cmidrule(lr){9-11}
        & & & & 操作系统 & 数据库 & 知识图谱 & 数字卡牌 & 横向谜题 & 家务整理 & 网页购物 & 网页浏览 \\
        \midrule
        \multirow{10}{*}{API}
        & \code{gpt-4}                & 0613 & 4.01 & 42.4 & 32.0 & 58.8 & 74.5 & 16.6 & 78.0 & 61.1 & 29.0 \\
        & \code{claude-3}             & opus & 3.11 & 22.9 & 51.7 & 34.6 & 44.5 & 14.3 & 70.0 & 27.9 & 26.0 \\
        & \code{glm-4}                & ---  & 2.89 & 29.2 & 42.3 & 46.3 & 34.1 & 14.2 & 34.0 & 61.6 & 27.0 \\
        & \code{claude-2}             & ---  & 2.49 & 18.1 & 27.3 & 41.3 & 55.5 &  8.4 & 54.0 & 61.4 &  0.0 \\
        & \code{claude}               & v1.3 & 2.44 &  9.7 & 22.0 & 38.9 & 40.9 &  8.2 & 58.0 & 55.7 & 25.0 \\
        & \code{gpt-3.5-turbo}        & 0613 & 2.32 & 32.6 & 36.7 & 25.9 & 33.7 & 10.5 & 16.0 & 64.1 & 20.0 \\
        & \code{text-davinci-003}    & ---  & 1.71 & 20.1 & 16.3 & 34.9 &  3.0 &  7.1 & 20.0 & 61.7 & 26.0 \\
        & \code{claude-instant}       & v1.1 & 1.60 & 16.7 & 18.0 & 20.8 &  5.9 & 12.6 & 30.0 & 49.7 &  4.0 \\
        & \code{chat-bison-001}      & ---  & 1.39 &  9.7 & 19.7 & 23.0 & 16.6 &  4.4 & 18.0 & 60.5 & 12.0 \\
        & \code{text-davinci-002}    & ---  & 1.25 &  8.3 & 16.7 & 41.5 & 11.8 &  0.5 & 16.0 & 56.3 &  9.0 \\
        \midrule
        \multirow{3}{*}{OSS（大）}
        & \code{llama-2-70b}          & chat & 0.78 &  9.7 & 13.0 &  8.0 & 21.3 &  0.0 &  2.0 &  5.6 & 19.0 \\
        & \code{guanaco-65b}          & ---  & 0.54 &  8.3 & 14.7 &  1.9 &  0.1 &  1.5 & 12.0 &  0.9 & 10.0 \\
        & \code{codellama-34b}        & instr.& 0.96 &  2.8 & 14.0 & 23.5 & 8.4 & 0.7 &  4.0 & 52.1 & 20.0 \\
        \multirow{4}{*}{OSS（中）}
        & \code{vicuna-33b}           & v1.3 & 0.73 & 15.3 & 11.0 &  1.2 & 16.3 &  1.0 &  6.0 & 23.9 &  7.0 \\
        & \code{wizardlm-30b}         & v1.0 & 0.46 & 13.9 & 12.7 &  2.9 &  0.3 &  1.8 &  6.0 &  4.4 &  1.0 \\
        & \code{guanaco-33b}          & ---  & 0.39 & 11.1 &  9.3 &  3.2 &  0.3 &  0.0 &  6.0 &  6.2 &  5.0 \\
        & \code{vicuna-13b}           & v1.5 & 0.93 & 10.4 &  6.7 &  9.4 &  0.1 &  8.0 &  8.0 & 41.7 & 12.0 \\
        & \code{llama-2-13b}          & chat & 0.77 &  4.2 & 11.7 &  3.6 & 26.4 &  0.0 &  6.0 & 25.3 & 13.0 \\
        & \code{openchat-13b}         & v3.2 & 0.70 & 15.3 & 12.3 &  5.5 &  0.1 &  0.0 &  0.0 & 46.9 & 15.0 \\
        & \code{wizardlm-13b}         & v1.2 & 0.66 &  9.0 & 12.7 &  1.7 &  1.9 &  0.0 & 10.0 & 43.7 & 12.0 \\
        & \code{vicuna-7b}            & v1.5 & 0.56 &  9.7 &  8.7 &  2.5 &  0.3 &  6.4 &  0.0 &  2.2 &  9.0 \\
        & \code{codellama-13b}        & instr.& 0.56 &  3.5 &  9.7 & 10.4 &  0.0 &  0.0 &  0.0 & 43.8 & 14.0 \\
        \multirow{7}{*}{OSS（小）}
        & \code{codellama-7b}         & instr.& 0.50 &  4.9 & 12.7 &  8.2 &  0.0 &  0.0 &  2.0 & 25.2 & 12.0 \\
        & \code{koala-13b}            & ---  & 0.34 &  3.5 &  5.0 &  0.4 &  0.1 &  4.4 &  0.0 &  3.9 &  7.0 \\
        & \code{llama-2-7b}           & chat & 0.34 &  4.2 &  8.0 &  2.1 &  6.9 &  0.0 &  0.0 & 11.6 &  7.0 \\
        & \code{codegeex2-6b}         & ---  & 0.27 &  1.4 &  0.0 &  4.8 &  0.3 &  0.0 &  0.0 & 20.9 & 11.0 \\
        & \code{dolly-12b}            & v2   & 0.14 &  0.0 &  0.0 &  0.0 &  0.1 &  1.2 &  0.0 &  0.4 &  9.0 \\
        & \code{chatglm-6b}           & v1.1 & 0.11 &  4.9 &  0.3 &  0.0 &  0.0 &  0.0 &  0.0 &  0.5 &  4.9 \\
        & \code{oasst-12b}            & sft-4& 0.03 &  1.4 &  0.0 &  0.0 &  0.0 &  0.0 &  0.0 &  0.3 &  1.0 \\
        \bottomrule
    \end{tabular}
\end{table*}

然而，本文所测试的 OSS LLM 通常在某些具有挑战性的任务——如知识图谱、
数字卡牌游戏、家务整理——上难以取得突破。我们在图~\ref{fig:oss-size} 中
描绘了它们的性能随模型规模的变化趋势。整体而言，绝大多数 OSS LLM 的
表现远逊于 API 类 LLM（平均得分 0.51 对 2.15）。在评测范围内（$\leq$70B）
表现最佳的 OSS LLM 为 \code{codellama-34b}，其总体得分为 0.96，但与
\code{gpt-3.5-turbo} 之间仍存在明显的性能差距。社区仍需付出大量努力，
以产出更强的 OSS LLM。

\subsection{分析}

在评测中，我们分析了若干影响 LLM 智能体在 AgentBench 上表现的关键因素，
包括执行结果的归因分析、代码训练的影响，以及 API 类商用 LLM 与 OSS LLM
之间的差异。关于规划、自我纠正与工具使用等方面的更多洞见与案例分析，
请参见附录~\ref{app:detailed}。

\paragraph{不同执行结果类型的占比。}
表~\ref{tab:outcomes} 报告了各类执行结果（参见第~\ref{sec:definition}~节）
的占比。AgentBench 任务中的主导失败原因为"超出任务限制"（TLE），这揭示了
LLM 在推理与决策能力上的薄弱。该数据为现有 LLM 的短板提供了佐证，
有望对未来研发提供指引（详见本文所提出的工具包）。

在数据库与数字卡牌游戏任务中，由于任务对输出格式的要求非常严格，
"格式无效"（IF）类错误较为常见（详见附录~\ref{app:detailed}）。相反，
在家务整理与网页浏览任务中，由于 LLM 经常生成超出预定义动作空间的动作，
"动作无效"（IA）类错误更为频繁。各类别在不同模型上的更详细占比，请参见
附录~\ref{app:detailed}。

\begin{table}[t]
    \centering
    \small
    \caption{
        \textbf{表 4：8 项任务中各类执行结果的占比（跨模型平均）}。
        CLE：超出上下文长度（Context Limit Exceeded）；
        TLE：超出任务限制（Task Limit Exceeded）。
        完整数据详见第~\ref{sec:definition}~节。
    }
    \label{tab:outcomes}
    \begin{tabular}{l c c c c c c c c}
        \toprule
        & 操作系统 & 数据库 & 知识图谱 & 数字卡牌 & 横向谜题 & 家务整理 & 网页购物 & 网页浏览 \\
        \midrule
        完成          & 75.0 & 37.9 & 30.1 & 51.2 & 14.0 & 13.1 & 54.9 & 56.6 \\
        CLE          &  0.1 &  0.7 &  2.0 &  0.0 &  3.5 &  0.7 &  0.0 &  0.0 \\
        格式无效      &  0.0 & 53.3 &  0.0 & 38.5 &  0.0 &  0.0 & 17.2 &  0.0 \\
        动作无效      &  0.9 &  0.0 &  0.0 & 10.2 &  0.0 & 64.1 &  0.0 &  8.4 \\
        TLE          & 23.9 &  8.0 & 67.9 &  0.0 & 82.5 & 22.1 & 27.8 & 35.0 \\
        \bottomrule
    \end{tabular}
\end{table}

\paragraph{代码训练的两面性影响。}
我们发现，代码微调会深度影响模型的推理生成与思维方式，其影响远不止于
"代码"这一议题本身。从 \code{codellama} 与 \code{llama-2} 系列模型的对比
可以看出：经过代码训练的模型在相对程式化、流程较为固定的任务
（例如网页购物）上表现更佳。然而，代码训练也可能会损害模型的通用思考
能力——例如 \code{codellama} 系列在数字卡牌游戏上的表现不及 \code{llama-2}
系列，在操作系统任务上（其中与 Linux 系统交互的重要性高于撰写 bash 代码）
亦如此。这一现象揭示了在 LLM 训练中，遵循流程的能力与通用思考能力之间的
权衡。

\paragraph{高质量对齐数据训练的影响。}
另一项富有启发的对比来自 \code{vicuna-13b} 与 \code{llama-2-13b}。
两者共享同一基座 LLM，但 \code{vicuna-13b} 是通过对 ShareGPT 数据
（由 \code{gpt-4} 与 \code{gpt-3.5-turbo} 生成、由用户共享）进行训练
完成对齐的，而 \code{llama-2-13b} 则是从零开始对齐。实验结果显示，
\code{vicuna-13b} 在 AgentBench 上的表现优于 \code{llama-2-13b}，甚至
与规模为其 3 倍的 \code{codellama-34b} 不相上下。这表明高质量对齐仍是
开发更优 LLM 智能体的关键。

\paragraph{意料之外的 \code{llama-2-13b} 与 \code{llama-2-70b} 表现相近现象。}
在实验中，我们意外地发现 \code{llama-2-13b} 与 \code{llama-2-70b}
的表现相近，尽管二者规模相差悬殊。在仔细复核并重跑实验后，结果依然如故。
我们认为，这反映了 \code{llama-2-70b} 的预训练不够充分。两者均使用 2T
token 进行预训练，但根据 \emph{缩放法则}（scaling law）
\citep{hoffmann2022}，更大的 LLM 应当使用更多 token 进行训练。

另一可能的原因是 \code{llama-2-70b} 在指令遵循对齐方面不充分，导致其在
指令遵循能力上相对滞后（参见附录~\ref{app:detailed}）。
